\begin{pastexam}[label={px:ntust-111-et-ds-05}]{臺科大電子所甲組 111 學年度 第 5 題}
（追蹤、問答）

(15\%) Consider the sequence $L$: 5, 55, 4, 44, 3, 33, 2, 22, 1, 11. Answer the following question:
\begin{enumerate}[label=(\arabic*)]
\item Write the status of the sequence $L$ after each phase of the insertion sort. (5\%)
\item Write the status of the sequence $L$ after each phase of the radix sort with radix $=10$. (5\%)
\item What is stable sorting? Please give a brief example to explain. (5\%)
\end{enumerate}
\end{pastexam}

\begin{pastsoln}
(1)(3) 屬資結書第~14~章（插入排序在第~1.3~節）；(2) 屬第~15~章。

\textbf{(1) 插入排序}（資結書虛擬碼~1.3）：第 $i$ 趟（$i = 2..10$）把第 $i$ 個數插進前面已經排好的 $i - 1$ 個數裡，
比它大的往後挪一格。 一共 $9$ 趟，每趟結束時的序列（直線左邊是已排好的部分）：
\begin{center}
\begin{tabular}{cll}
\toprule
趟 & 插入 & 結束時的 $L$ \\
\midrule
1 & 55 & 5, 55 $\mid$ 4, 44, 3, 33, 2, 22, 1, 11 \\
2 & 4 & 4, 5, 55 $\mid$ 44, 3, 33, 2, 22, 1, 11 \\
3 & 44 & 4, 5, 44, 55 $\mid$ 3, 33, 2, 22, 1, 11 \\
4 & 3 & 3, 4, 5, 44, 55 $\mid$ 33, 2, 22, 1, 11 \\
5 & 33 & 3, 4, 5, 33, 44, 55 $\mid$ 2, 22, 1, 11 \\
6 & 2 & 2, 3, 4, 5, 33, 44, 55 $\mid$ 22, 1, 11 \\
7 & 22 & 2, 3, 4, 5, 22, 33, 44, 55 $\mid$ 1, 11 \\
8 & 1 & 1, 2, 3, 4, 5, 22, 33, 44, 55 $\mid$ 11 \\
9 & 11 & 1, 2, 3, 4, 5, 11, 22, 33, 44, 55 \\
\bottomrule
\end{tabular}
\end{center}
第 1 趟 $55 > 5$，不用挪。 這組輸入每插入一個較小的個位數，前面所有的數都要往後挪，挪動次數就是逆序數對的個數
（$0 + 2 + 1 + 4 + 2 + 6 + 3 + 8 + 4 = 30$ 次）。

\textbf{(2) 基數排序}（LSD，資結書第~15.2.2~節）：基數 $10$，最大的數是兩位數，做兩趟，每趟依一個位數把數依序
丟進 $0$--$9$ 號桶子（同一個桶子保持原本的先後順序），再依桶子編號收回來。
\begin{itemize}
\item 第 1 趟，依個位數：桶 1：$1, 11$；桶 2：$2, 22$；桶 3：$3, 33$；桶 4：$4, 44$；桶 5：$5, 55$。 收回得
\[ 1, 11, 2, 22, 3, 33, 4, 44, 5, 55 . \]
\item 第 2 趟，依十位數（一位數的十位是 $0$）：桶 0：$1, 2, 3, 4, 5$；桶 1：$11$；桶 2：$22$；桶 3：$33$；桶 4：$44$；
桶 5：$55$。 收回得
\[ 1, 2, 3, 4, 5, 11, 22, 33, 44, 55 . \]
\end{itemize}
第 2 趟裡 $1, 2, 3, 4, 5$ 在桶 0 的順序，是第 1 趟排出來的順序；每趟都穩定，個位數的排序結果才會留下來。

\textbf{(3) 穩定排序}：鍵相同的元素，排序後的相對順序和排序前相同（資結書第~14.1~節）。 例：依成績排序
$\langle (\text{甲}, 80), (\text{乙}, 70), (\text{丙}, 80) \rangle$。 甲、丙都是 $80$，排序前甲在丙前面。 穩定的排序
由小到大得 $\langle (\text{乙}, 70), (\text{甲}, 80), (\text{丙}, 80) \rangle$，甲仍在丙前面；不穩定的排序可能得到
$\langle (\text{乙}, 70), (\text{丙}, 80), (\text{甲}, 80) \rangle$。 用途：資料原本依姓名排好，再用穩定排序依成績排，
同分的人仍然依姓名排列；(2) 的基數排序也正是靠每一趟穩定才正確。 插入排序、合併排序、氣泡排序、計數排序是
穩定的；選擇排序、堆積排序、快速排序一般不穩定。
\end{pastsoln}
